\section{Variabilidad entre Máquinas}

\subsection{Naturaleza del Problema}

La variabilidad entre máquinas no constituye una excepción puntual, sino una característica estructural del sistema a construir. Cada máquina puede presentar diferencias en los datos que envía, en los protocolos de comunicación que utiliza, en la frecuencia de actualización, en la lógica de cálculo y en los procesos operativos asociados.

Por tanto, el diseño del sistema debe asumir desde el inicio que no existe un modelo único y cerrado de máquina aplicable a todos los casos.

\subsection{Dimensiones de Variabilidad}

La variabilidad identificada afecta, al menos, a las siguientes dimensiones:

\begin{itemize}
    \item \textbf{Fuentes de datos}: una máquina puede integrarse mediante MQTT, HTTP u otros mecanismos.
    \item \textbf{Parámetros}: cada máquina puede exponer campos distintos, con diferente semántica, unidades y tipos de dato.
    \item \textbf{Transformaciones}: los datos recibidos pueden requerir mapeos, conversiones o normalizaciones específicas.
    \item \textbf{Reglas de negocio}: una máquina puede necesitar reglas de cálculo o interpretación propias.
    \item \textbf{Alertas}: las condiciones que constituyen una incidencia pueden variar entre máquinas.
    \item \textbf{Operación}: algunas máquinas pueden trabajar con órdenes, pausas, contadores o procesos manuales, mientras que otras no.
\end{itemize}

\subsection{Estrategia de Tratamiento}

La estrategia propuesta consiste en mantener un núcleo común de plataforma y desplazar la variabilidad hacia definiciones, configuraciones y servicios especializados.

En consecuencia:

\begin{itemize}
    \item La plataforma compartirá un mismo modelo base para máquinas, lecturas, estados y alertas.
    \item Cada máquina podrá disponer de sus propios parámetros definidos explícitamente.
    \item Las reglas y métricas específicas se describirán mediante definiciones propias.
    \item La lógica compleja o delicada podrá implementarse en servicios específicos por máquina o por familia de máquinas.
\end{itemize}

\subsection{Evitar Dos Riesgos Comunes}

Este enfoque también pretende evitar dos riesgos frecuentes en soluciones de este tipo.

El primer riesgo consiste en modelar todas las máquinas como si fueran iguales, introduciendo posteriormente excepciones y condiciones especiales repartidas por toda la aplicación. Este enfoque genera rigidez y un mantenimiento costoso.

El segundo riesgo consiste en intentar resolver toda la variabilidad a través de un motor completamente genérico desde la primera iteración. Aunque atractivo en apariencia, este enfoque suele introducir una complejidad excesiva antes de disponer de suficiente conocimiento real del dominio.

La solución recomendada se sitúa en un punto intermedio: una base común estable, combinada con mecanismos de extensión progresivos y controlados.

\subsection{Criterio de Escalabilidad}

El sistema se considerará correctamente escalable si la incorporación de una nueva máquina implica principalmente:

\begin{itemize}
    \item registrar la máquina en la plataforma,
    \item definir sus parámetros y fuentes de entrada,
    \item configurar o implementar sus reglas particulares,
    \item y exponer sus datos y estado sin alterar el núcleo común.
\end{itemize}

Este criterio debe guiar las siguientes fases del diseño y de la implementación.